\section{Proof of the Iteration Lemma}\label{sec:pf_it}
We give a proof of the iteration lemma. Let $\val_f(H)$ denote the fraction of edges inside $H$ that are satisfied by the assignment $f$.

\begin{lemma}\label{lem:subroutine-restated}
Let $c,\gamma,\delta \in (0,1]$ and $r \in \nn$. Let $I$ be an affine UG instance on alphabet $\Sigma$ on $J_{n,\ell,\alpha}$ with $\ell, n$ large enough  and value at least $c$. 
%Let $\delta_{\min} = \min_{\eta \in [c-\gamma,c]}(\delta(\eta))$.
Suppose we have a subroutine $\cA$ which given as input any affine UG instance $I'$ on $J_{n,\ell,\alpha}$ with $\val_\mu(I') \ge c-\gamma$, returns an $\leq r$-restricted subcube $H$ on $J$ and a partial assignment $f$ such that, $\text{val}_{f}(H) \geq \delta$.
% = \delta(c-\val_\mu(I))
Then if $\cA$ runs in time $T(\cA)$, there is a $O(|V(J)|T(\cA)+|V(J)|^3)$-time algorithm which finds a solution for $I$ that satisfies an $\Omega(\delta\gamma(1-\alpha)^r)$-fraction of the  edges of $J$. 
\end{lemma}

\begin{proof}
We will use the algorithm $\cA$ as a subroutine. To get a full assignment, our algorithm below is a generalized version of the Algorithm~\ref{alg:j}, where we've replaced some of the steps in Algorithm~\ref{alg:j} with an arbitrary subroutine $\cA$ that finds a subcube and an assignment to it with high value. We include it here for completeness.

\begin{algorithm-thm}[Partial to Full Assignment]\label{alg:partial} ~\\ 
\begin{compactenum}
\item Set $j = 1, I_0 = I, R_0 = \phi$.
\item While $|R_{j-1}| < \frac{\gamma}{2}|V(J)|$:
\begin{compactenum}
\item Run subroutine $\cA$ on $I_{j-1}$ to find an $\leq r$-restricted subcube $H_{j}$ and partial assignment $f_j$.
\item Let $S_j$ be the induced subgraph of $H_j$ induced by the set $H_j \setminus R_{j-1}$ and assign them using $f_j$. Let $R_j = R_{j-1} \cup H_j$. 
\item Randomize the constraints on edges incident on $R_j$ and let the new instance by $I_j$. 
\item Increment $j$.
\end{compactenum}
\item Output any assignment to $V(G)$ that agrees with all partial assignments $f_j$ (assigned to $S_j$'s) considered above.
\end{compactenum}
\end{algorithm-thm}


Let us analyse the $j^{th}$-iteration of this procedure. We are given an instance $I_{j-1}$.

\begin{claim}\label{claim:it-drop}
$\val(I_{j-1}) \geq \val(I) - \frac{2|V(R_{j-1})|}{|V(G)|}$.
\end{claim}

\begin{proof}
Suppose $X$ is an assignment that achieves the value of $I$. We can check that $X$ satisfies at least $\val(I) - \frac{2|V(R_{j-1})|}{|V(G)|}$-fraction of edges of $I_{j-1}$, hence the claim follows. To see this note that the fraction of edges incident on $R_{j-1}$ and therefore the fraction that is randomized is at most $\frac{2|V(R_{j-1})|}{|V(G)|}$ and every edge not incident on $R_{j-1}$ remains satisfied by $X$.
\end{proof}

We will show that the partial assignment $f_j$ to $H_j \setminus R_{j-1}$ satisfies a large fraction of edges (as counted in the whole graph). Before that let us show that the randomization of edges ensures that the value of the edges randomized is small (even when restricted to some subcube). Let $E(S)$ denote the edges incident on $S$ (including those with only one endpoint in $S$) and let $E(S,S)$ denote the edges with both endpoints in $S$. For a set of edges $E$ let $\val_f(E)$ denote the fraction of edges in $E$ that are satisfied by $f$.


\begin{claim}\label{claim:chernoff}
With probability $1-o(1)$ the instance $I_{j-1}$ satisfies the following for all $r$-restricted subcubes $C$:
\[\val(E(C \cap R_{j-1})) \leq \frac{2}{|\Sigma|},\]
where $\val(E(S))$ denotes the maximum fraction of $E(S)$ that can be satisfied by any UG assignment to the endpoints of $E(S)$.
\end{claim}

\begin{proof}
Let $k$ denote $|\Sigma|$. Fix an $i$-restricted subcube $C$ for $i \leq r$. We know that $R_{j-1}$ is a union of $\leq r$-restricted subcubes, therefore $C\cap R_{j-1}$ also forms a union of $\leq r$-restricted subcubes inside $C$. Let this collection of subcubes be denoted by $\cC$.
Consider the set of edges $E$ inside $C$ on which the constraints were randomized (in the previous iteration), and let $A$ be the set of vertices in $C$ that are endpoints of them. We claim that with probability $1-o(1)$,  for all assignments $F$ to $A$, $F$ satisfies at most $2/k$ of the constraints. First, the number of assignments is $k^{|A|}$, so we fix one assignment $F$, analyze it, and then union bound over all assignments.


Fix $F$; for each $e\in E$, the probability that $F$ satisfies $e$ is $1/k$. Hence in expectation $F$ satisfies
$\mu = 1/k$ fraction of the constraints, and by Chernoff’s bound
\[\Pr[\val_F(E) \geq 2\mu ]\leq e^{-\frac{|E|}{3k}}.\]

Now we lower bound $|E|$. Note that if $v\in A$, then either $v$ was in a subcube in $\cC$ or $v$ was adjacent to a vertex in a subcube from $\cC$. Let the degree of $v$ inside $C$ be denoted by $d$ which we know is ${\ell-i \choose \alpha(\ell - i)}{n-\ell \choose \alpha(\ell-i)}$.

\begin{enumerate}
    \item If $v \in C'$ for some $C' \in \cC$ then at least $(1-\alpha)^{r}$ of the edges incident to $v$ remain in $C'$ and therefore are in $E$. 
    \item If $v$ is adjacent to $u \in R_{j-1}\cap C$ then $u \in C'$ for some $C' \in \cC$. In particular there is a subset $I \subseteq u$ of size at most $r$ such that $C'$ is an $I$-restricted subcube. It follows that in fact $v$ has at least ${\ell-i \choose \alpha(\ell-i)}{n-\ell-r \choose \alpha(\ell-i) – r}$ neighbours in $C'$. By a simple calculation this is at least $n^{-r} d$, and it follows that at least $n^{-r} d$ of the edges adjacent to $v$ are in $E$.
\end{enumerate}

In conclusion, in any case we get that at least $n^{-r}d$ of the edges adjacent to $v \in A$ are inside $E$, hence $|E| \geq n^{-r} d |A|/2 \gg 3n|\Sigma| |A|$, so Chernoff gives a probability of $e^{-n|\Sigma| |A|}$ which is good enough for
a union bound over all ($k^{|A|}$ many) assignments to $A$. Additionally we union bound over all $\leq r$-restricted subcubes $C$ which are only $n^{O(r)}$ many, to get the statement of the lemma.
\end{proof}

At iteration $j$ of the while-loop, since the while condition is met, we know that $I_{j-1}$ has value $\geq c - \gamma$ (Claim~\ref{claim:it-drop}) so inside the while-loop, $\cA$ will always find an $\leq r$-subcube $H_j$ and assignment $f_j$ with $\val_{f_j}(H_j) \geq \delta$. Next, we find an assignment $f_j$ to the set of vertices $V(S_j)$ that by definition don't intersect previously assigned vertices. Since $f_j$ doesn't reassign any vertices, in the final step of the algorithm it is possible to output an assignment that is consistent with all previously considered partial assignments. We will now show that our final partial assignment satisfies a large fraction of the edges, where we say that an edge $(u,v)$ is satisfied by a partial assignment $f_j$, if both vertices $u,v$ have been assigned labels under $f_j$ and the labels satisfy the edge.


\begin{claim}\label{claim:it-val}
The value of the partial assignment found at iteration $j$ satisfies:
\[\val(f_j) \geq \delta(1 - \alpha)^r\frac{|V(H_j)|}{2|V(G)|},\]
where $\val(f_j)$ denotes the fraction of edges (in $E(G)$) satisfied by the partial assignment $f_j$.
\end{claim}

\begin{proof}
Let $H_j = S_j \cup (H_j \cap R_{j-1})$ and let $E_1 = E(S_j,S_j)$ and $E_2 = E(H_j \cap R_{j-1})$ so that $E_1 \cup E_2 = E(H_j,H_j)$ denoted by $E$. By Claim~\ref{claim:chernoff} we know that $\val_{f_j}(E_2) \leq 2/k$ where $k$ denotes $|\Sigma|$. 
By an averaging argument:
\[\delta \leq \frac{\#\text{sat}_{f_j}(E_1)}{|E|}+ \frac{|E_2|}{|E|}\val_{f_j}(E_2) \leq \frac{\#\text{sat}_{f_j}(E_1)}{|E|} + \frac{2}{k},\]
where $\#\text{sat}_{f_j}(E_1)$ denotes the number of edges of $E_1$ satisfied by $f_j$. This implies that $\frac{\#\text{sat}_{f_j}(E_1)}{|E|} \geq \delta/2$. 

Since $H_j$ is a $\leq r$-restricted subcube at least $(1-\alpha)^r$-fraction of edges stay inside $H_j$, hence $|E|/|E(H_j)|$ is at least $(1-\alpha)^r$. Using that $|E(H_j)|/|E(J)| \geq |V(H_j)|/|V(J)|$ we get that: 
\[\val(f_j) =  \frac{\#\text{sat}_{f_j}(E_1)}{|E|} \cdot \frac{|E|}{|E(H_j)|} \cdot \frac{|E(H_j)|}{|E(J)|} \geq \frac{\delta}{2}(1 - \alpha)^r\frac{|V(H_j)|}{|V(G)|},\]
where $\val(f_j)$ denotes the edges satisfied by $f_j$ (assigned to $S_j$) as a fraction in the whole graph. 
\end{proof}

Once we have these facts, the conclusion is immediate. Firstly there cannot be more than $V(G)$ iterations of the while-loop, since at each iteration we assign at least one new vertex. Each iteration takes time $T(\cA) + |V(G)|^2$, hence the algorithm runs in time $|V(G)|(T(\cA) + |V(G)|^2)$. 

Suppose the algorithm exits the while loop at the $t+1$ iteration. Then we know that $R_{t} \geq \frac{\gamma}{2}|V(J)|$.
Using claims~\ref{claim:it-val} above, we get that the value of the final assignment is proportional to the number of vertices assigned: 
\[\sum_{j=1}^t\val(f_j) \geq \delta(1 - \alpha)^r \sum_{j=1}^t \frac{|V(H_j)|}{2|V(J)|} \geq \delta(1 - \alpha)^r\frac{|R_t|}{2|V(J)|} \geq \delta(1 - \alpha)^r\frac{\gamma}{4},\]
thus proving the lemma.
\end{proof}
\section{Missing Proofs}\label{sec:missing_pfs}

The following claim bounds the Booleanity error term 
and shows that there is a \sos proof that they are
negligible (provided that we take the parameter $\nu$ 
to be sufficiently small).
\begin{claim}\label{claim:bool}
\[\cA_I ~~ \vdash_{\poly(\ell^r/\nu)} ~~ \E_{c}[\sum_s B(H^c_s)] \leq |\Sigma|\ell^{O(r)}\sqrt{\nu}.\]
\end{claim}

\begin{proof}
Recall that $B(F) = \frac{4}{3}\E_u[(F(u)^3 - F(u))(\Pi_{\geq \lambda_r}F)(u)] = \frac{4}{3}\ip{F^3 - F, PF}$, where $\Pi$ is the projection operator $\Pi_{\geq \lambda_r}$. Using Cauchy-Schwarz we get that for all $\eta > 0$:
\begin{equation}\label{eq:bool-cs}
B(F) \leq \frac{4}{3}\left(\frac{\eta}{2}||F^3 - F||_2^2 + \frac{1}{2\eta}||\Pi F||_2^2\right) \leq \frac{2\eta}{3}\E_u[F(u) -F(u)^3] + \frac{2}{3\eta}\delta(F^2),
\end{equation}
where the last inequality holds since $F \in [0,1]$ and $\Pi$ is a projection operator. We will use this 
to bound $B(H_s^c)$ as follows. 
Fixing $s, u$, taking $F = H_s^c(u)$, 
the first term above is 
$\E_c[(H_s^c(u)-H_s^c(u)^3]$ and we upper bound it 
by $O(\ell^r\nu)$. Indeed, first as 
$0\leq H_s^c(u)\leq 1$ we have
\begin{equation}\label{eq:bool1}
\E_c[H_s^c(u) - H_s^c(u)^3] = (1+H_s^c(u))\E_c[H_s^c(u) - H_s^c(u)^2] \leq 2 \E_c[H_s^c(u) - H_s^c(u)^2].    
\end{equation}
Next, recall that $H_s^c(u)$ is a product of the terms $(1-R^c_{s,a})$ for $a \subset u, |a| \leq r$. Additionally each variable $R^c_{s,a}$ is a product of approximate indicators: namely there are $\beta_i \in (0,1)$  such that $R_{s,a}^c(u) = p^{c,s,a}_1(\delta^{c,s,a}_1)\cdot \ldots\cdot p^{c,s,a}_k(\delta^{c,s,a}_k)$, where $p^{c,s,a}_i = p_{\beta_i+c\nu^2,\nu^2}$ or $p_{<\beta_i+c\nu^2,\nu^2}$ as $c$ varies in $\{0,1,\ldots,\lfloor 1/\nu \rfloor\}$, and $\delta_i$'s are linear functions of the variables $G_s(u)$ and lie in $[0,1]$ under the axioms $\cA_I$.
By Claim~\ref{claim:silly_sos}(applied twice) it follows that:
\[\E_c[H_s^c(u) - H_s^c(u)^2] \leq \sum_{a \subset u: |a|\leq r}\E_c[R^c_{s,a} - (R^c_{s,a})^2] \leq \sum_{a \subset u: |a|\leq r}\sum_{i \in [k]}\E_c[p_i^{c,s,a}(\delta^{c,s,a}_i) - p_i^{c,s,a}(\delta^{c,s,a}_i)^2],\]
so we can bound each of the terms above.  

We will show that each term is smaller than $O(\nu)$. Let $p$ be an approximate indicator polynomial $p = p_{\beta,\nu}$ and let $p^c$ denote $p_{\beta+c\nu^2,\nu^2}$ for $c \in \{0,\ldots,\lfloor 1/\nu\rfloor\}$. For all $\delta \in [0,1]$, we argue that
\begin{equation}\label{eq:bool-random}
\delta \in [0,1] \vdash_{\tO(1/\nu^2)} ~~ \E_c[p^c(\delta) - p^c(\delta)^2] \leq O(\nu).
\end{equation}
Indeed, this follows by case analysis on $\delta \in [0,1]$ (that we show next), and then using L\'{u}kacs theorem (Corollary~\ref{cor:Lukacs}) we get there is an \sos proof of degree $\deg(p^c)$. 

If $\delta \in [0,\beta] \cup [\beta+\nu, 1]$ then one can check that $p^c(\delta) \in [0,\nu^2]$ for all $c$. Therefore let $\delta \in [\beta,\beta+\nu]$, and specifically in one of the intervals - $[\beta+t\nu^2, \beta+(t+1)\nu^2]$ for $t \in \{0,\ldots,\lfloor1/\nu\rfloor\}$. Then we get that $p^t(\delta)$ could be any number between $[0,1]$ and therefore $p^t(\delta)-p^t(\delta)^2 \leq 1/4$. But for every $t' \neq t$ we get that $p^{t'}(\delta)$ is either in the interval $[0,\nu^2]$ or $[1-\nu^2,1]$ implying that $p^{t'}(\delta) - p^{t'}(\delta)^2 \leq O(\nu^2)$. Averaging over $a$ we therefore get:
\[\E[p^c(\delta) - p^c(\delta)^2] \leq \nu \cdot \frac{1}{4}+ (1-\nu)O(\nu^2) \leq O(\nu).\]
In conclusion, we get that 
$\E_c[H_s^c(u) - H_s^c(u)^2]\leq O(k\ell^r\nu) = \ell^{O(r)}\nu$. Plugging this into~\eqref{eq:bool-cs} 
and taking $\eta = 1/\sqrt{\nu}$ we get that
\[\E_c[B(H^c_s)] \leq  \sqrt{\nu}\ell^{O(r)} + O(\sqrt{\nu}\E_c[\delta(H^c_s)]) \leq \sqrt{\nu}\ell^{O(r)},\]
and summing over $s \in \Sigma$ finishes the proof.
\end{proof}

\begin{claim}\label{claim:silly_sos}
\[
\{0 \leq x_i \leq 1\} \vdash_{2k} 
\prod_{i \in [k]}x_i - (\prod_{i \in [k]}x_i)^2 \leq x_1-x_1^2 + \ldots + x_k - x_k^2. 
\]
\end{claim}
\begin{proof}
We prove this by induction. For $k=1$ this is clear, and for $k=2$ we have
\begin{align*}
xy - (xy)^2 &= xy – xy^2 + xy^2 – x^2 y^2 \\
&= x(y-y^2)+ (x-x^2) y^2 \\
&\leq y-y^2 + x-x^2.    
\end{align*}

For the inductive step, let $k\geq 3$ and denote $y=x_1\cdots x_{k-1}$. Applying the base case and 
then the inductive hypothesis we get
\[\prod_{i \in [k]}x_i - \left(\prod_{i \in [k]}x_i\right)^2 
=y x_k - (y x_k)^2
\leq y-y^2 + x_k-x_k^2
\leq x_1 - x_1^2 + \ldots + x_{k-1}^2 - x_{k-1}^2.
\]
One can check that the SoS degree used is at most $2k$.
\end{proof}



% The following lemma asserts several properties of shift symmetric pseudodistributions.


% \begin{lemma}[Shift-Symmetry properties]\label{lem:sym}
% Let $\mu$ be a degree-$D$ shift-symmetric pseudodistribution satisfying the unique games axioms $\A_{I}$ given by (\ref{eq:ip}) for an affine Unique-Games instance $I$.
% Then:
% \begin{enumerate}
%     \item For all vertices $u,v \in V(G)$ and shifts $s \in \Sigma$,
%     $\pPr[X_v = s \mid X_u = 0] = \pPr[X_v - X_u = s]$.
%     \item For all polynomials $f(X)$ with $\deg(f) \le D-2$, such that $f(X) = f(X+s)$ for every global shift $s \in \Sigma$,
%     \[\pE[f(X) \mid X_v-X_u=s] = \pE[f(X) \mid X_u = 0, X_v = s].\]
% \end{enumerate}
% \end{lemma}

% \begin{proof}
% Recall that since $\mu$ is a shift-symmetric pseudodistribution, we have that,
% \begin{equation}\label{eq:props}
% \pE_{\mu}[X_{u_1,a_1} \cdots X_{u_m,a_m}]
% = \frac{1}{|\Sigma|}\sum_{t}\pE_\mu[X_{u_1,a_1 - t}\cdots X_{u_m,a_m -t}] = \pE_{\mu}[X_{u_1,a_1+s} \cdots X_{u_m,a_m+s}]
% \end{equation}
% for all $\{(u_1,a_1), \ldots, (u_m,a_m)\} \in ([n] \times \Sigma)^{\le D}$ and $s \in \Sigma$. 

% For the first item, note that in $\mu$ the marginal distribution of each $X_v$ is uniform, so $\pPr_{\mu}[X_v = s]=1/|\Sigma|$ for all $s \in \Sigma$ and so 
% \begin{align*}
% \pPr_{\mu}[X_v = s \mid X_u = 0]
% &= \frac{\pPr_{\mu}[X_v = s, X_u = 0]}{\pPr_{\mu}[X_u = 0]}\\
% &= |\Sigma| \cdot \pPr[X_v = s, X_u = 0] \\
% &= \sum_{t \in \Sigma} \pPr[X_v = s+t, X_u = t]\\
% &= \pPr[X_v - X_u = s].
% \end{align*}

% For the second item, fix $f$ to be any polynomial as in the second item, and note that
% \begin{align*}
% \pE_{\mu}[f(X) \mid X_v - X_u = s]
% &= \frac{\pE_{\mu}[\sum_{t \in \Sigma} f(X) \cdot \Ind[X_v = s + t, X_u = t]]}{\pE_{\mu}[\sum_{t \in \Sigma}\Ind[X_v = s+t, X_u = t]]}\\
% &= \frac{\pE_{\mu}[\sum_{t \in \Sigma} f(X + t) \cdot \Ind[X_v = s+t, X_u = t]]}{\pE_{\mu}[\sum_{t \in \Sigma}\Ind[X_v = s+t, X_u = t]]}\\
% &= \frac{|\Sigma|\cdot\pE_{\mu}[f(X) \cdot \Ind[X_v = s, X_u = 0]]}{|\Sigma|\cdot\pE_{\mu}[\Ind[X_v = s, X_u = 0]]}\\
% &= \pE_{\mu}[f(X) \mid X_v = s, X_u = 0],
% \end{align*}
% where to obtain the second equality we have used the shift-symmetry of $f$, $f(X) = f(X+(t-s))$, and in the penultimate equality we have used the shift-invariance of $\mu$.
% \end{proof}




% It follows along the lines of~\cite{BBKSS} but we give it here for completeness. Before doing that let us define the Condition\&Round value of a subgraph.

% \begin{definition}[Condition\&Round Value]\label{def:cr-val}
% Given a unique games instance $I = (G,\Pi)$ and a degree $4$ shift-symmetric pseudodistribution $\mu$ over $I$, for every subgraph $H$ of $G$, let $\text{ind-val}_\mu(H)$ denote the expected fraction of satisfied edges when independent rounding is performed on $V(H)$ using the marginals of $\mu$, i.e. $\text{ind-val}_\mu(H) := \E_{(v,w) \sim E(H)}[\sum_s \pE_\mu[X_{v,s}]\pE_\mu[X_{w,\pi_{vw}(s)}]]$. Let the Condition\&Round value, denoted by $\text{CR-val}_\mu(H)$ be the value obtained by performing Algorithm~\ref{alg:low-ent-restated} on $H$, i.e. $\text{CR-val}_{\mu}(H) := \E_{u \sim V(H)}[\text{ind-val}_{\mu \mid X_u = 0}(H)]$.
% \end{definition}



% \begin{proof}
% We will use the algorithm $\cA$ as a subroutine. To get a full assignment, our algorithm below is a generalized version of the Algorithm~\ref{alg:j}, where we've replaced the steps 2 to 4 in Algorithm~\ref{alg:j} with an arbitrary subroutine $\cA$ that finds a subgraph with high Condition\&Round value with respect to $I$. We include it here for completeness.

% \begin{algorithm-thm}[Partial to Full Assignment]\label{alg:partial} ~\\ 
% \begin{compactenum}
% \item Solve the degree-$D$ SoS SDP relaxation for the integer program $\cA_I$ and make the pseudodistribution shift-symmetric to get a pseudodistribution $\mu_0$. Set $j = 1$.

% \item While $\eta_{j-1} := c - \val_{\mu_{j-1}}(I) \leq \gamma$:
% \begin{compactenum}
% \item Run subroutine $\cA$ on $\mu_{j-1}$ to find a subgraph $H_{j}$ with $\text{SubRound-val}_{\mu_{j-1}}(H_j) \geq \delta(\eta_{j-1})$.
% \item Run SubRound (Algorithm~\ref{alg:j-partial}) on $H_j$ with pseudodistribution $\mu_{j-1}$ to get an assignment $f_j$ to $V(H_j)$. Let $S_j$ be a subgraph of $H_j$ induced by the set of vertices that have not been previously assigned by any partial assignment $f_{k}, k < j$ and assign them using $f_j$. 
% \item Rerandomize the pseudodistribution $\mu_{j-1}$ on $S_j$ to get $\mu_j$: Make the marginal distribution over the assigned vertices uniform and independent of other vertices, that is, for all degree $\leq D$ monomials define
% $\pE_{\mu_{j}}$ as follows, 
% \[\pE_{\mu_{j}}[X_{h_1,a_1} \cdots X_{h_t,a_t}X_{u_1,b_1} \cdots X_{u_m,b_m}]
% := \frac{1}{|\Sigma|^t}\pE_{\mu_{j-1}}[X_{u_1,b_1}\cdots X_{u_m,b_m}],
% \]
% where $\{(h_1,a_1), \ldots, (h_t,a_t)\} \in (V(S_j) \times \Sigma)^{t}$ and $\{(u_1,b_1), \ldots, (u_m,b_m)\} \in (({[n] \choose \ell} \setminus V(S_j)) \times \Sigma)^{m}$.
% \item Increment $j$.
% \end{compactenum}
% \item Output any assignment to $V(G)$ that agrees with all partial assignments $f_j$ considered above.
% \end{compactenum}
% \end{algorithm-thm}


% Let us first check that the algorithm is well-defined. The initial pseudodistribution $\mu_0$ by definition satisfies axioms $\cA_I$ and is shift-symmetric. In subsequent iterations of the while loop all these properties are satisfied by the modified pseudodistributions: 1) the rerandomizing operation on pseudodistributions produces a valid pseudodistribution operator that satisfies the axioms $\cA_I$ and is also shift-symmetric, 2) At iteration $j$ of the while-loop, since the while condition is met, we know that $\mu_{j-1}$ has value $\geq c - \gamma$ and furthermore we can show that since the value only decreases at each step, it is always $\leq c$, so that $\eta_j \in [0,\gamma]$. So inside the while-loop, $\cA$ will always find a non-empty subgraph $H_j$ with high Condition\&Round value. Next, we find an assignment $f_j$ to the set of vertices $V(S_j)$ that by definition don't intersect previously assigned vertices. Since $f_j$ doesn't reassign any vertices, in the final step of the algorithm it is possible to output an assignment that is consistent with all previously considered partial assignments. We will now show that our final partial assignment satisfies a large fraction of the edges, where we say that an edge $(u,v)$ is satisfied by a partial assignment $f_j$, if both vertices $u,v$ have been assigned labels under $f_j$ and the labels satisfy the edge. We claim the following two facts:

% \begin{claim}
% The drop in value in every iteration satisfies that:
% \[\val_{\mu_{j-1}}(I) - \val_{\mu_j}(I) \leq \frac{2|V(H_j)|}{|V(G)|},\]

% where $\val_\mu(I)$ denotes the SDP value of $I$ with respect to the pseudodistribution $\mu$. 
% \end{claim}

% \begin{proof}
% For any edge $(h,v)$ where $h \in H_j$, we have that $\val_{\mu_j}((h,v)) = \frac{1}{k}$, whereas $\val_{\mu_{j-1}}((h,v)) \leq 1$. For any edge whose both endpoints lie outside $H_j$, the value remains unchanged under rerandomizing. Noting that the fraction of edges incident on vertices in $H_j$ is at most $\frac{2|V(H_j)|}{|V(G)|}$ the conclusion follows.
% \end{proof}


% \begin{claim}
% The value of the partial assignment found at iteration $j$ satisfies:
% \[\val(f_j) \geq \delta^2(\eta_{j-1})(1 - \phi)\frac{|V(H_j)|}{2|V(G)|},\]
% where $\val(f_j)$ denotes the fraction of edges (in $E(G)$) satisfied by the partial assignment $f_j$.
% \end{claim}

% \begin{proof}
% Since the SubRound value of $H_j$ is large, we know that there exists a pseudodistribution $\cD$ on $X$ that is obtained by conditioning $\mu_{j-1}$ on an event $Y_A = y_a$ (for $|A| = O(1)$) such that $\text{CR-val}_{\cD^\sym}(H_j) \geq \delta(\eta_{j-1})$. We will first prove that $\text{CR-val}_{\cD^\sym}(S_j) \geq \delta(\eta_{j-1}) := \delta$, where $S_j$ is the subgraph induced by the unassigned (by previous partial assignments $f_k$, $k < j$) vertices of $H_j$. For notational simplicity we will drop the subscript $j$ from $H_j,S_j$. 


% We know by the guarantees of the subroutine $\cA$ that $H$ is such that, $\text{CR-val}_{\cD^\sym}(H) = \E_{u \sim V(H)}[\text{ind-val}_{\cD^\sym | X_u = 0}(H)] \geq \delta$.  First note that the marginals of every vertex are uniform, due to the shift-symmetry of $\cD^\sym$. Moreover we have that conditioning on previously assigned vertices, i.e. any vertex $u \in H \setminus S$, maintains this property, since the distribution of $u$ is completely uncorrelated with the other vertices. So we get that, $\text{ind-val}_{\mu | X_u = 0}(H) = \frac{1}{|\Sigma|} < \delta$ for all $u \in H \setminus S$. This implies that, 
% \begin{equation}\label{eq:cr1}
% \E_{u \sim S}[\text{ind-val}_{\mu | X_u = 0}(H)] \geq \text{CR-val}_{\mu}(H) \geq \delta.    
% \end{equation}
% Again we have that, $\text{ind-val}_{\mu | X_u = 0}(e) = \frac{1}{|\Sigma|}$, for any edge $e$ which has at least one endpoint in $H \setminus S$, so we get that, $\E_{u \sim S}[\text{ind-val}_{\mu | X_u = 0}(S)] \geq \delta$. Now we can perform derandomized Condition\&Round on $S$ to get an assignment $f_j$ that satisfies at least a $\delta$-fraction of the edges of $S$.

% Next we will show, by an averaging argument, that the edges of $S$ constitute a large fraction of the edges incident on the vertices of $H$, which would imply that $f_j$ satisfies a large fraction of these edges. Let $u_0 \in S$ be a vertex for which $\text{ind-val}_{\mu | X_{u_0} = 0}(H) \geq \delta$ (we know such a vertex exists by equation~(\ref{eq:cr1})) and let $\mu'$ be the pseudodistribution $(\mu | X_{u_0} = 0)$. First note that the set of edges $E(H) \setminus E(S)$, have independent rounding value $1/|\Sigma|$ under $\mu'$, since at least one endpoint of such edges has been assigned previously. Since $\text{ind-val}_{\mu'}(H) \geq \delta$, a simple averaging argument gives us that the set $E(H) \setminus E(S)$ can be at most a $\frac{1 - \delta}{1 - (1/|\Sigma|)}$-fraction of $E(H)$. So the set $E(S)$ is at least a $\frac{\delta - (1/|\Sigma|)}{1 - (1/|\Sigma|)} \geq \delta/2$-fraction of $E(H)$. Since the expansion of $H$ is at most $O(\eta_{j-1})$, we have that $E(H)$ is a $(1 - O(\eta_{j-1}))$-fraction of the total edges incident on $H$, which in turn is at least a $\frac{|V(H)|}{|V(G)|}$-fraction of $E(G)$. Combining these facts we get that $f_j$ satisfies a $\delta \cdot \frac{\delta}{2} \cdot (1 - O(\eta_{j-1}))\cdot \frac{|V(H)|}{|V(G)|}$-fraction of the edges of $G$. 
% \end{proof}

% Once we have these facts, the conclusion is immediate. Firstly there cannot be more than $V(G)$ iterations of the while-loop, since  at each iteration we set the value of at least one new vertex to $1/|\Sigma|$. The rerandomization operation in the while loop as well as the symmetrization operation (Lemma~\ref{lem:sym}) can be done in time polynomial in the description of $\mu_0$. So each iteration takes time $T(\cA) + |V(G)|^{O(D)}$, hence the algorithm runs in time $|V(G)|(T(\cA) + |V(G)|^{O(D)})$. 

% Moreover, combining the claims above, we get that the value of the partial assignment at any iteration is proportional to the drop in value of the pseudodistribution. That is, 

% \[\val(f_j) \geq \delta^2(\eta_{j-1})(1 - O(\eta_{j-1}))\frac{|V(H_j)|}{2|V(G)|} \geq \delta^2(\eta_{j-1})(1 - O(\eta_{j-1}))\left(\frac{\val_{\mu_{j-1}}(I) - \val_{\mu_j}(I)}{4}\right).\]

% At the last iteration, we know that the pseudodistribution value has dropped by at least $\gamma$ (compared to $\val_{\mu_0}(I)$), hence summing the above over all iterations $j$, we get that the value of the partial assignment returned by the algorithm is at least ${\frac{1}{4}(\min_j(\delta(\eta_j)))^2(1 - O(\eps))\eps}$ as required.
% \end{proof}









